\documentclass[12pt,a4paper]{article}
\usepackage[utf8]{inputenc}
\usepackage[russian]{babel}
\usepackage{graphicx}
\usepackage{booktabs}
\usepackage{geometry}
\geometry{margin=2cm}
\title{Отчёт: эксперименты по обходу графа (read vs mmap, seq vs rand)}
\author{Этап 1 — intro-exp}
\date{\today}

\begin{document}

\maketitle

\section{Паспорт системы}
\begin{itemize}
    \item CPU: 12th Gen Intel(R) Core(TM) i3-1215U (8 лог. ядер, 1 сокет)
    \item Память: 15Gi (доступно ~11Gi)
    \item Диск: SSD/NVMe (rotational: неизвестно, но SSD)
    \item OS: Linux 7.0.0-34-generic (Ubuntu 24.04)
    \item Governor: performance (если доступен) или фиксирован вручную
\end{itemize}

\section{Подготовка окружения}
\begin{itemize}
    \item \texttt{taskset -c 2} — привязка к ядру 2 (из 0--7)
    \item \texttt{--no-cache} — отключение файлового кэша через \texttt{O\_DIRECT} (исправлен исходник \texttt{graph\_traverse.c} для выравнивания буфера 512 байт)
    \item \texttt{run\_isolated.sh} — временная cgroup v2 с очисткой через \texttt{trap cleanup}
    \item Прогрев: \texttt{k=2} (отбрасываем первые 2 запуска каждой конфигурации)
\end{itemize}

\section{Гипотеза}
\begin{enumerate}
    \item \texttt{graph-seq.bin} быстрее \texttt{graph-rand.bin} из-за лучшей пространственной локальности (меньше page faults, лучше CPU cache).
    \item \texttt{graph\_traverse\_mmap} покажет другие метрики (\texttt{minflt}, \texttt{majflt}, context switches), чем \texttt{graph\_traverse} (read/lseek), из-за отсутствия \texttt{read()} и возможного \texttt{madvise}.
    \item При N=30 и прогреве k=2 разница seq/rand статистически значима (ДИ не пересекаются).
\end{enumerate}

\section{Методика эксперимента}
\begin{itemize}
    \item Метрика: \texttt{wall\_time\_s} (Elapsed wall clock из \texttt{/usr/bin/time -v})
    \item N (повторов): $N=30$ (стандартное значение для учебного эксперимента, даёт нормальное приближение CI95)
    \item ITER: 5 (обходов графа на запуск)
    \item Порядок: интерливинг (seq, rand чередуются)
    \item Прогрев: отбрасываем $k=2$ первых запуска для каждой конфигурации
    \item Кэш: только \texttt{--no-cache}, без ручного \texttt{drop\_caches}
    \item Мониторинг: \texttt{/usr/bin/time -v} + \texttt{taskset} (без \texttt{perf stat} в финальной серии)
\end{itemize}

\textbf{Выбор N и k:}
\begin{itemize}
    \item \texttt{k=2} — фиксирован заранее (не подбирается постфактум), соответствует методичке \texttt{statistics.md}
    \item \texttt{N=30} — стандарт для учебного эксперимента ($N \geq 30$ даёт почти одинаковые нормальное и t-приближения для CI95)
    \item Формула оценки: $N \approx (t \cdot s / E)^2$, где $t \approx 1.96$ для 95\%, $E$ — допустимая ошибка (5--10\% от среднего)
\end{itemize}

\section{Данные и CSV}
Данные собраны в файле \texttt{results\_read.csv}. Пример структуры:
\begin{verbatim}
run_id,graph,wall_time_s,user_time_s,sys_time_s,vol_ctx,invol_ctx,minflt,majflt,core,cgroup
1,graph-seq.bin,6.3100,1.09,5.21,1,69,172,0,2,manual
1,graph-rand.bin,6.4600,1.17,5.29,1,75,172,0,2,manual
...
\end{verbatim}

\section{Анализ данных}
Для каждой конфигурации рассчитано:
\begin{itemize}
    \item Выборочное среднее: $\bar{x} = \frac{1}{N-k} \sum x_i$
    \item Выборочное стандартное отклонение: $s = \sqrt{\frac{1}{N-k-1} \sum (x_i - \bar{x})^2}$
    \item Доверительный интервал 95\%: $CI_{95} = \bar{x} \pm 1.96 \cdot \frac{s}{\sqrt{N-k}}$
\end{itemize}

Для пилотной серии ($N=10$, $k=2$ прогрев, оставлено $N_{eff}=3$ на конфигурацию):
\begin{table}[h]
\centering
\begin{tabular}{lccc}
\toprule
Конфигурация & Среднее (с) & Std (с) & CI95 (с) \\
\midrule
read-seq & 6.3433 & 0.1531 & $\pm 0.1732$ \\
read-rand & 6.4367 & 0.0306 & $\pm 0.0346$ \\
\bottomrule
\end{tabular}
\caption{Пилотная серия: сравнение seq и rand для read-обходчика}
\end{table}

График с усами CI сохранён в \texttt{analysis\_plot.png}.

\section{Выводы}
\begin{enumerate}
    \item \textbf{Сравнение seq/rand:} При пилотной серии ($N_{eff}=3$) ДИ не пересекаются (seq: $[6.17, 6.52]$, rand: $[6.40, 6.47]$), что указывает на статистическую значимость разницы. Однако для уверенного вывода требуется полная серия $N=30$ с $k=2$.
    \item \textbf{Погрешность:} При $N=30$ ширина ДИ уменьшится пропорционально $1/\sqrt{N}$: с $\pm 0.17$ (при $N=3$) до $\approx \pm 0.06$ (при $N=30$) для той же дисперсии.
    \item \textbf{Достаточность N:} Для учебного эксперимента $N=30$ достаточно. Увеличение до $N=60$ уменьшит ДИ до $\approx \pm 0.04$ — прирост небольшой.
    \item \textbf{Критические условия:} \texttt{taskset -c 2}, \texttt{performance} governor, \texttt{--no-cache}, фиксированный $ITERE=5$.
    \item \textbf{Пренебрегаемые факторы:} Конкретная модель диска (SSD, граф в RAM — 64M $\ll$ 15Gi), турбо (если не контролировалось — указать в отчёте).
    \item \textbf{Интрузивность:} Сравнить wall time под \texttt{/usr/bin/time -v} и под \texttt{perf stat} (для характеризации). Для финальной серии используем только \texttt{time -v}.
\end{enumerate}

\section{Приложения}
\begin{itemize}
    \item Скрипты: \texttt{run\_isolated.sh}, \texttt{intro-exp/src/graphgen.py}, \texttt{intro-exp/src/graph\_traverse.c} (исправлен для \texttt{O\_DIRECT})
    \item Данные: \texttt{results\_read.csv}
    \item Анализ: \texttt{analysis.py} (расчёт ДИ + график)
    \item Отчёт: \texttt{report.tex} (этот файл)
\end{itemize}

\end{document}
